\chapter{Installazione}
\label{ch:installazione}

Guida cross-platform per installare ed eseguire piTantum su Windows,
Linux e macOS. La versione \emph{di riferimento} di questa guida \`e
\texttt{docs/installation.md} nel repo (Markdown, sempre aggiornato);
qui se ne riproduce il contenuto in forma compatta per il manuale
PDF. In caso di divergenza, fa fede il file Markdown.

\paragraph{Architettura del lancio.} Il progetto \`e composto da due
processi:
\begin{itemize}
  \item \textbf{Backend}: \texttt{uvicorn} che serve l'app FastAPI in
        \texttt{webui/backend} sulla porta 8000.
  \item \textbf{Frontend}: server di sviluppo \texttt{vite} che serve
        l'app SvelteKit in \texttt{webui/frontend} sulla porta 5173.
\end{itemize}
Uno script unico (\texttt{start.bat} su Windows,
\texttt{start.sh} su Linux/macOS) li avvia entrambi. Al primo lancio
crea la \emph{venv} Python, installa i requirements, fa
\texttt{npm install}; ai lanci successivi parte in pochi secondi.

% ----------------------------------------------------------------
\section{Windows}

\paragraph{Prerequisiti.}
\begin{itemize}
  \item Python $\geq$ 3.11 -- installer ufficiale:
        \url{https://www.python.org/downloads/windows/}
  \item Node.js $\geq$ 20 LTS -- installer ufficiale:
        \url{https://nodejs.org/en/download} (scegliere ``LTS'').
  \item Git -- installer ufficiale:
        \url{https://git-scm.com/download/win}.
\end{itemize}

\paragraph{Note importanti.}
\begin{itemize}
  \item Sull'installer Python, \textbf{spuntare} ``Add python.exe to
        PATH''. Senza, lo script non trova il comando \texttt{python}.
  \item Riavviare il terminale dopo le installazioni: il PATH non
        si aggiorna nelle finestre g\`ia aperte.
\end{itemize}

\paragraph{Clone + lancio.}
\begin{lstlisting}[style=shell]
git clone https://github.com/gborghi/timetable.git
cd timetable
webui\start.bat
\end{lstlisting}
Lo script apre due finestre cmd (``piTantum -- backend'' e
``piTantum -- frontend''). Aprire il browser su
\url{http://localhost:5173}.

\paragraph{Stop.} Chiudere le due finestre cmd, oppure
\texttt{Ctrl+C} in ciascuna.

\paragraph{Troubleshooting Windows.}
\begin{itemize}
  \item ``\texttt{python} non riconosciuto'': l'installer Python non
        \`e stato eseguito con ``Add to PATH''. Reinstallare oppure
        aggiungere \texttt{C:\textbackslash Users\textbackslash
        \emph{tu}\textbackslash AppData\textbackslash Local\textbackslash
        Programs\textbackslash Python\textbackslash Python311} al PATH.
  \item Porta 8000 occupata: \texttt{netstat -ano | findstr :8000}
        + \texttt{taskkill /PID <pid> /F}.
  \item Antivirus blocca \texttt{start.bat}: aggiungere alle
        eccezioni di Windows Defender.
\end{itemize}

% ----------------------------------------------------------------
\section{Linux (Ubuntu, Debian, Fedora, Arch)}

\paragraph{Prerequisiti.} Python $\geq$ 3.11, Node $\geq$ 20 LTS,
Git, build-essential.

\paragraph{Ubuntu / Debian.}
\begin{lstlisting}[style=shell]
sudo apt update
sudo apt install -y python3 python3-venv python3-pip nodejs npm \
                    git build-essential
\end{lstlisting}

I repo \texttt{apt} di Ubuntu 22.04 / Debian 12 hanno Node \emph{12-18},
troppo vecchio per SvelteKit 2 (richiede 20+). In quel caso usare
\textbf{nvm}:
\begin{lstlisting}[style=shell]
curl -o- https://raw.githubusercontent.com/nvm-sh/nvm/v0.39.7/install.sh | bash
source ~/.bashrc
nvm install --lts && nvm use --lts
\end{lstlisting}

\paragraph{Fedora.}
\begin{lstlisting}[style=shell]
sudo dnf install -y python3 python3-virtualenv nodejs npm git \
                    gcc make
\end{lstlisting}

\paragraph{Arch / Manjaro.}
\begin{lstlisting}[style=shell]
sudo pacman -S python python-pip nodejs npm git base-devel
\end{lstlisting}

\paragraph{Clone + lancio.}
\begin{lstlisting}[style=shell]
git clone https://github.com/gborghi/timetable.git
cd timetable
chmod +x webui/start.sh webui/stop.sh
./webui/start.sh                # background; log in webui/logs/
./webui/start.sh --foreground   # log live in console (Ctrl+C ferma)
\end{lstlisting}

\paragraph{Stop.}
\begin{lstlisting}[style=shell]
./webui/stop.sh
\end{lstlisting}

\paragraph{Troubleshooting Linux.}
\begin{itemize}
  \item Porta in uso: \texttt{sudo lsof -i :8000} + \texttt{kill <pid>}
        (eventualmente \texttt{kill -9 <pid>}).
  \item \texttt{python3} mancante su distro minimal: usare
        l'eseguibile chiamato \texttt{python} (3.x).
  \item WSL2: \texttt{start.sh} funziona normalmente. Aprire il
        browser dal lato Windows su
        \url{http://localhost:5173}; WSL2 inoltra le porte
        automaticamente.
\end{itemize}

% ----------------------------------------------------------------
\section{macOS (Intel + Apple Silicon)}

\paragraph{Prerequisiti.}
\begin{itemize}
  \item Xcode Command Line Tools: \texttt{xcode-select --install}.
  \item Homebrew: vedere \url{https://brew.sh/}.
  \item Python $\geq$ 3.11, Node $\geq$ 20 LTS, Git: tutti via Homebrew.
\end{itemize}

\paragraph{Installazione.}
\begin{lstlisting}[style=shell]
# Homebrew (se non presente)
/bin/bash -c "$(curl -fsSL https://raw.githubusercontent.com/Homebrew/install/HEAD/install.sh)"

# Toolchain
brew install python@3.12 node git

# Verifica
python3 --version    # >= 3.11
node --version       # >= v20
\end{lstlisting}

\paragraph{Note Apple Silicon.}
\begin{itemize}
  \item \texttt{ortools} dalla v9.7 spedisce wheel ARM64 nativi:
        Rosetta non serve.
  \item Se Python \`e installato da python.org invece che Homebrew,
        scegliere l'installer ``macOS 64-bit universal2''.
\end{itemize}

\paragraph{Clone + lancio + stop.} Identico a Linux.

\paragraph{Troubleshooting macOS.}
\begin{itemize}
  \item Gatekeeper blocca \texttt{start.sh}:
        \texttt{xattr -d com.apple.quarantine webui/start.sh
        webui/stop.sh}.
  \item \texttt{ortools} senza wheel ARM (versioni vecchie):
        \texttt{pip install --upgrade ortools} dentro la venv del
        backend.
  \item \texttt{brew} non in PATH:
        \texttt{eval "\$(/opt/homebrew/bin/brew shellenv)"} in
        \texttt{\textasciitilde/.zprofile}.
\end{itemize}

% ----------------------------------------------------------------
\section{Verifica installazione}

Dopo il primo lancio andato a buon fine:

\begin{lstlisting}[style=shell]
python --version       # >= 3.11   (Linux/macOS: spesso python3)
node --version         # >= v20
npm --version          # >= 10

# Backend health
curl http://127.0.0.1:8000/api/health
# atteso: {"status":"ok","name":"pitantum","version":"0.1.0"}

# Async layer (Section 2.5 P2)
curl http://127.0.0.1:8000/api/health/async
# atteso: {"status":"ok","async":true,"dialect":"sqlite"}

# Frontend
curl -I http://127.0.0.1:5173
# atteso: HTTP/1.1 200 OK
\end{lstlisting}

% ----------------------------------------------------------------
\section{Aggiornamento}

Quando si fa \texttt{git pull} per portare a casa nuovi commit:

\begin{lstlisting}[style=shell]
cd /path/to/timetable
git pull origin main

# se sono cambiate dipendenze Python:
webui/backend/.venv/bin/pip install -r webui/backend/requirements.txt
# (Windows: webui\backend\.venv\Scripts\pip install ...)

# se sono cambiate dipendenze frontend:
cd webui/frontend && npm install && cd ../..

# se ci sono migration DB:
cd webui/backend && ./.venv/bin/alembic upgrade head && cd ../..
# (Windows: webui\backend\.venv\Scripts\alembic upgrade head)
\end{lstlisting}

In pratica \texttt{start.sh}/\texttt{start.bat} rilanciato dopo un
\texttt{git pull} gestisce automaticamente \texttt{pip install} e
\texttt{npm install} quando servono. Le migration DB sono coperte da
\texttt{\_apply\_lightweight\_migrations()} allo startup come
safety-net (vedi capitolo \emph{Modello dati}); per la via canonica
usare \texttt{alembic upgrade head}.

Dopo un aggiornamento backend, fare un \textbf{hard reload} del
browser (\texttt{Ctrl+Shift+R}) per pulire chunk JavaScript cached.

% ----------------------------------------------------------------
\section{Disinstallazione}

\begin{lstlisting}[style=shell]
cd /path/to/timetable
./webui/stop.sh                 # Linux/macOS
# (Windows: chiudere le finestre cmd di start.bat)

rm -rf webui/backend/.venv      # ~150 MB
rm -rf webui/frontend/node_modules  # ~400 MB
rm -rf webui/data/timetable.db  # DATI UTENTE (salvarli prima!)
\end{lstlisting}

Per rimuovere completamente, cancellare la cartella \texttt{timetable/}.
Python, Node, Git restano installati nel sistema; rimuoverli dal
gestore di pacchetti solo se non servono ad altri progetti.

% =========== Cap 3: Architettura ===========
